\documentclass[12pt,a4paper]{article}
\usepackage[scheme=plain]{ctex}
\usepackage[margin=2.5cm]{geometry}
\usepackage{indentfirst}
\usepackage{setspace}
\usepackage{fontspec}
\setCJKmainfont[BoldFont=SimHei,ItalicFont=KaiTi]{SimSun}
\newCJKfontfamily\kaifont{KaiTi}[AutoFakeBold]
\usepackage{newpxtext,newpxmath}

\usepackage{lastpage}





\pagestyle{plain}

\usepackage[nobottomtitles*]{titlesec}
\titleformat{\section}{\Large\bfseries}{}{0em}{ }
\titlespacing*{\section}{0pt}{3ex}{1.5ex}

\usepackage{xcolor}
\usepackage{needspace}

\setlength{\parindent}{2em}
\setlength{\parskip}{0.4em}

\doublespacing

\title{\bfseries Day 7\\[0.3em]收束与回望}
\author{}
\date{}

\begin{document}
\maketitle

\vspace{-2em}

\vspace{1em}

\section{聘用始末}
\vspace{-0.3em}{\color{gray}\small\kaifont \noindent K 不只要一个简单的说法。他想搞清楚，聘他来做土地测量员的这个决定是如何做出的。村长讲了很久。故事的核心不是答案。}\vspace{0.5em}
\vspace{0.3em}

\noindent\rule{\textwidth}{0.4pt}

\vspace{0.3em}



“这么说这两个助手，”村长又开口了，洋洋得意地微笑着，似乎在说：所有这些全是他一手安排的，而谁都无法预想到这一点，“这么说，他们让您感到讨厌，可他们是您自己的助手呵。”——“不对，”K冷冷地说，“他们是我到此地后才跑到我身边来的。”——“怎么，跑来的，”村长说，“您的意思大概是说派来的吧。”——“好吧，就算是派来的吧，”K说。“但说他们是从天外飞来的也未尝不可，这种派法也够欠考虑的了。”——“这里办什么事都不会欠考虑，”村长说，急得连腿疼也忘了，一骨碌坐起来。“办什么事都不欠考虑，”K说，“那么招聘我这事又怎么说？”——“聘用您也是经过深思熟虑的，”村长说，“只是有一些次要因素起了干扰作用罢了，我可以用有关的文件向您证明这一点。”——“这些文件恐怕找不到了，”K说。“找不到？”村长叫起来。“米齐，请快点找！不过，没有文件我也可以先把事情的经过讲给您听听。刚才我提到的那份指令，我们接到后当即回复上头表示感谢，说我们不需要土地测量员。但是这封回信看来是没有送回发出指令的那个部去，我管它叫A部吧，而是误送到另一个部，送到B部去了。这样A部就一直没有得到回音，然而可惜连B部也没有收到我们复信的全部；一种可能是公文袋里装的信落在村公所，再一种可能是在半路上遗失了——肯定不会在部里遗失，这我敢担保，总之，B部也只收到一个空公文袋，袋上只注明封套内所装的、可惜实际上并不在内的文件是关于聘任土地测量员一事的。A部呢，一直还在等着我们的回复，这个部虽然已将此事记录在案，但部主管人指望着我们会答复，等着在收到答复之后他们或者聘任土地测量员，或者根据需要继续同我们就此事进行磋商。这种坐等下面回信的事绝非绝无仅有，这是可以理解的，在办什么事都讲求准确、一丝不苟的地方这样做也无可厚非。可是结果呢，因为他指望着我们，所以也就不大重视部里那份有关记录，弄到后来就把这事忘记了。但是在B部就不同，公文袋送到了一位以认真负责闻名的部主管人手里，他叫索尔蒂尼，是个意大利人；连我这个内行人也不明白为什么像他那样能干的人会被长期安置在那个可说很低的职位上。这位索尔蒂尼自然把空文件袋派人送回给我们要求补上内容。但是，从A部发出第一封函件到那时已经过了好几个月，甚至已经好几年了；这也不难理解，因为按常规如果一份文件送对路了，那么最迟在第二天也就送到部里，当天就可以处理；但如果没送对路——在组织非常严密完善的情况下，文件要走这条错误的路简直就必须费很大力气去找，否则是找不到的，那么，那么当然就会拖很长时间了。因此当我们收到索尔蒂尼的便函时，就只能隐隐约约地记起这件事来，当时我们管这事的只有两人，就是米齐和我，那位老师还没有分配给我，我们两个只能对最重要的事保存一份副本，简短地说吧，我们只能给他一个很笼统的回答，说我们完全不知道聘任一事，我们这里也不需要什么土地测量员。”



“但是，”村长顿住了，似乎发觉在乐此不疲的叙述中忘其所以，把话扯远了，或者至少有可能扯远了，“这个故事您听着不觉得腻味吧？”



“不，”K说，“我觉得挺有意思，挺解闷的。”



然后村长说：“我讲这些给您听不是为了给您解闷。”



“我所以觉得它有意思，”K说，“仅仅因为我听了这事后，对那种可笑的混乱算是窥见了一斑，这种混乱有时能决定一个人的命运哟。”



“您还没有真正窥见一斑，”村长正色说，“我可以接着跟您讲。我们的回答当然满足不了索尔蒂尼那样的人。我很佩服这个人，尽管一听到他的名字就头疼。为什么？因为他谁都不信赖，比如，即使他有过无数次机会进行了解，完全知道那是个最可靠的人，可到下一次再碰上时他还是不信任那个人，好像他根本不认识人家似的，或者更正确些说，好像他知道那人是个无赖。我觉得这是对的，一个公职人员就应该这样；可惜我的天性使我不能按这条原则办事，你不是看见了吗，我现在对您这个外乡人把什么底都亮出来了，我这叫做禀性难移，没法子。而索尔蒂尼一接到我们的回信就起了疑。于是便有了一连串的通信往来。索尔蒂尼问我为什么突然想起说不要聘用土地测量员；我仗着米齐的好记性回信说，这事最初是上头提出来的呀（至于事实上是另一个部发来的文件，这一点我们早忘记了）；索尔蒂尼对此的说法是：为什么我到现在才提起上级的这封公函；我又回复他说：因为我现在才想起这封公函来嘛；索尔蒂尼呢：这真是太奇怪了；我：对于拖了那么长时间的一件事，这是一点也不奇怪的；索尔蒂尼：这事确实很奇怪，因为我想起来的那封公函并不存在；我：当然不存在啦，因为关于这事的全部文件都丢失了；索尔蒂尼：如果确有那第一封函件，就必定会有一条有关的记录在，然而这样一条记录并不存在。这时我无话可说了，因为要说是索尔蒂尼那个部出了差错，这话我是既不敢讲也不相信的。您，土地测量员先生，也许会暗暗埋怨索尔蒂尼，心想只要他考虑一下我的话，至少会去向另外的一些部打听打听这事吧。然而真要这样做恰恰是不对的，我不想让您哪怕只是在心里产生一种索尔蒂尼身上还有点毛病的印象。压根不考虑有出错的可能性，正是衙门的一条工作原则。由于组织十分完备严密，这条原则是正确的，而如果想达到极高的办事效率，它也是必要的。所以说索尔蒂尼根本就不能去向其他各部打听，而且就算他去打听，这些部门也决不会答复他，因为它们马上就会感觉出：这是在追究一个出差错的可能性。”



“村长先生，请允许我打断您一下，向您提个问题，”K说，〔8〕“您不是提到过一个检查机关吗？按照您对这里办事情况的描述，如果设想一下这里对这些事居然会没有监督检查，这简直令人不堪忍受。”



“您非常严格，”村长说。“但您就是再严格一千倍，同衙门对待自己的那种严格相比，您的严格也仍然等于零。只有一个地地道道的外乡人才会提出您提的问题。要问有没有检查机关吗？这里有的只是检查机关。当然啰，这些检查机关不是为了查出一般意义上的错误而设立的，因为这里是不会出错的，即便某次出了一个错，如您的情况，但是谁敢下定论，说这肯定是一个错误！”



“这个说法真是太新鲜了！”K叫起来。



“对我来说这是很旧很旧的说法，”村长说。“我同您差不多，也确信是出了一个错，索尔蒂尼因为对这事伤透脑筋毫无进展而得了重病，第一批检查机关官员来了，亏得他们揭出了错误的来源，看出这里有一个错，可是谁敢拍着胸脯说，第二批检查机关官员也会做出同样的判断，第三批和以后各批也都会得出同样的结论？”



“就算是这样吧，”K说，“不过我最好还是不要对这些考虑妄加评论，我也是头一次听说这些检查机关，当然还不可能完全了解它们。但我觉得这里必须区别两点：第一是机关内部发生的事以及仍然从机关的角度出发对这些事可能有这样或那样的看法，第二是我作为一个具体的个人，我这个身在公职机关以外的人，现在正面临着公职机关无视我的存在的危险，这种无视现实是如此荒谬，以致我直到现在总不大相信这危险是实际存在着的。对于第一点，很可能您村长先生刚才讲的那些话都是对的，您对内情了解得如此了如指掌令人叹服，只是我在听了您这番话后也想听您说说我个人的事情了。”



“我也正想来谈这个，”村长说，“但如果我不预先说明几句，那您是听不懂我的话的，我刚才提起检查机关看来是为时过早了。所以现在我要再回过头来谈谈我跟索尔蒂尼的分歧。我已经说过，渐渐的我招架不住了。但索尔蒂尼哪怕只比他的对手略胜一筹，他就等于大获全胜了，因为这样一来他办起事来精力就更集中、精神就更饱满，就更加指挥若定；结果是他的手下败将一见他就胆战心惊，而他手下败将的敌人看着他就会拍案叫绝。只因为我对这后一点在另外一些事情上也有亲身体会，我才能像现在这样讲他这个人。不过我还从未有幸亲眼见过他，他不能到下面来，他真是日理万机，公务繁忙至极，有人对我描述过他的房间，说四壁全被大捆大捆的卷宗挡满，一捆摞一捆，形成了一根根高大的方柱，而且这些还仅仅是索尔蒂尼当时正在处理的文件，由于不断从一捆捆卷宗中抽出、插进文件，并且又都是为赶时间匆匆忙忙地做，这些卷宗堆成的柱子就不断倒塌下来，正是这种接二连三、每隔一会儿就出现一次的轰然巨响，成了索尔蒂尼办公室的突出特征。这也难怪，索尔蒂尼是个一丝不苟的工作人员，不论是最小的事还是最大的事，他都一视同仁，一样认真对待。”



“村长先生，”K说，“您总把我的事叫做小事一桩，但我这件事使许多公职人员忙个不亦乐乎，或许它在开始时只是一桩很小的事吧，可是由于像索尔蒂尼那样的官员们积极参与，它确实已经变成一件大事了。这一点我感到非常遗憾，这是完全违反我的本意的，我不敢有那样的野心：想让关于我的事情的卷宗也变成高大的柱子又随时哗啦啦塌下来，我只希望做个小小的土地测量员，在一张小小的绘图桌旁安安心心地工作就心满意足了。”



“不，”村长说，“您的事不是大事，在这点上您没有理由抱怨，您的事确实是诸多小事中最小的一件。工作量大小并不决定工作的重要性，如果您这样想那您对衙门的了解就太皮毛了。但即使工作量大小能说明问题是否重要，您这件事也仍然是一粒微不足道的芝麻，普遍的事由，即那些没有出所谓差错的事，办理起来工作量就大得多，工作的成果自然也丰富得多，不过您还根本不知道您的事给人造成的实际上的工作量呢，我现在就来说说这个。首先要说的是，索尔蒂尼虽然不再追问我，但他手下的工作人员来了，每天都要在贵宾楼对一些有声望的村民进行几次质询并作翔实记录。大部分人支持我的说法，然而也有少数人起了疑心；土地测量问题触及农民的根本利益，农民们嗅出有人在算计他们，觉得这里头有些不公道的事，他们还找到了一个领头的，于是索尔蒂尼从他们提供的情况中得出结论，确信：如果我当初把问题提到村民大会上讨论，是不会出现全体一致反对聘用土地测量员的情况的。〔9〕这样一来一件原本是明摆着谁都明白的小事——即不需要土地测量员——不管怎么说至少被弄得成为一个问题了。在这件事上，有一个叫布伦斯维克的表现得特别突出——您大概不认识他吧，他也许不是坏人，但又蠢又狂，他是拉塞曼的一个妹夫。”



“是那个鞣皮匠拉塞曼吗？”K问，接着便描述了他在拉塞曼家见到的那个大胡子的长相。



“是的，正是他，”村长说。



“我还认识他的女人，”K有点有一搭没一搭地说。



“这是有可能的，”村长说，然后便沉默了。



“她很美，就是气色不大好，病恹恹的。她大概是城堡里来的人吧？”这后一句话是带着疑问语气说出来的。



村长看了看表，将药水倒进一把汤匙里，急急忙忙灌下肚去。



“您大概只认得城堡里那些办事机构吧？”K颇不客气地问道。



“对，”村长带着自嘲然而却是感激的微笑说道。“机关也是最重要的呵。至于说到布伦斯维克：要是我们能把他从村里驱逐出去，几乎对每个人来说都是件大快人心的事，拉塞曼也不会是最不开心的一个。可是当时布伦斯维克给自己拉了一些支持者：他虽然不是演说家，倒也会大喊大叫，而这样对某些人也能起作用。于是，终归我还是被迫把这件事提交村民大会讨论。不过这也就是布伦斯维克取得的唯一的胜利了，因为到了村民大会上，当然是绝大多数人反对聘用土地测量员的。这次大会也又过去好几年了，但是这些年来这件事一直没有平息，这部分是因为索尔蒂尼太认真，他通过极为周密细致的调查，力图摸清多数派和反对派的动机，再就是因为布伦斯维克的愚蠢和野心，他同官府有各种各样的私人关系，而且有层出不穷的、异想天开的新招数去利用这些关系。当然索尔蒂尼不会上布伦斯维克的当，布伦斯维克怎么瞒得了索尔蒂尼？——可是，恰恰是为了不致上当受骗，需要进行新的调查，新的调查还没有结束，布伦斯维克又想出了新花招，他这人很灵活，这正是他那愚蠢的一种表现。好，现在我来谈谈我们政府机构的一个特殊性质。与它的精细严密相适应，它也极为敏感。如果一件事已经深思熟虑了很久很久，那么即使考虑还没有结束也会出现这种情况：在某个事先无法预见事后也无法找到的地方会闪电般突然冒出一个解决办法，这办法虽说多半是正确地了结了那件事，但无论如何总带有一定的随意性吧。情形就好像是：政府机关再也受不了那绷得太紧的弦、受不了同一件也许本身是微不足道的事情那长年累月的刺激，从而自发地、不等哪个官员来过问就作出决定了。这当然不是说出现了什么奇迹，可以肯定地说是某位官员挥笔写下了这个解决方案，或是作了一个口头的决定，但至少从我们这里，从此地，甚至从公职机关也怎么都不可能弄清楚究竟是哪位官员在这事上作了定夺，是根据什么理由这样决定的了。这些全是检查机关在很久以后才查明的；而我们呢，就永远不得而知了，实际上这事过了那么久恐怕也不会再有谁感兴趣了。我说过，正是这类决定，它们多半是正确的，是很好的，只有一点不足，就是——这类事通常都是如此——一般人知道得太晚，因此在知道前的一段长时间里一直还在对那早已决定了的事进行热烈的讨论。我不知道在您的事情上是否也发出过一个这样的决定。——有迹象说明发出过，也有迹象说明没有发出过；假如已经发出，那么聘任书就可能早已寄给了您，然后您就会长途跋涉到此地来，这些花费了很多时间，而与此同时索尔蒂尼却还会一直在这里为这同一件事操劳着，累得筋疲力尽，布伦斯维克也还会绞尽脑汁搞阴谋诡计，而我就会在他们两人中间受尽了夹板气。这个可能性只是我的揣测，我确切知道的是这一点：一个检查机关查出来好多年前A部曾经就聘任土地测量员一事向村里发出过询问函件，而后直至现在一直没有收到回音。新近又有官员向我查询这事，这样一来全部事实当然也就真相大白了，A部对我提供的不需要土地测量员这一回答表示满意，索尔蒂尼则不得不看到：原来他并不主管这件事情，并且还做了许多无用的、让人伤透脑筋、把人弄得精疲力竭的虚功。——当然，他是无辜的。要不是这段时间也同以往一样新的工作任务不断从四面八方源源而来，要不是您的事情说到底仅仅是桩很小的事——几乎可以说是诸多小事当中最小的一件，那么我们大家一定都会感到如释重负了，我相信索尔蒂尼自己也会这样，只有布伦斯维克有一肚子怨气，但那只会让人觉得可笑罢了。好，现在，恰恰是现在，在整个这件事总算善始善终顺利了结之后——从那时起到今天又已经过去好久，您突然又冒了出来，事情似乎又像要从头开始的样子，请您，土地测量员先生，设想一下我有多么失望呵。现在我下了大决心在我力所能及的范围内决不让这类事再发生，这一点想来您是很清楚也能理解的吧？”
\par\vspace{0.3em}{\footnotesize\color{gray}（第五章）}
\newpage
\section{根茎}
\vspace{-0.3em}{\color{gray}\small\kaifont \noindent 德勒兹与加塔利说，世界上有三种生命空间形态。树根，统一的生命体，从唯一起点出发，通过二元分化不断生长。胚根，矛盾的生命体，多个起点相互嫁接，任何路径都可以被替代，整体在崩溃和重建之间轮回。根茎，多元的生命体，脱离起源，不断迁移、断裂，在任何一个节点上都蕴含着重新开始的潜力。}\vspace{0.5em}
\vspace{0.3em}

\noindent\rule{\textwidth}{0.4pt}

\vspace{0.3em}



书的第一种类型，是“书—根”（livre-racine）。树已经是世界的形象，或者说，根是“世界—树”的形象。这是经典之书，作为有机的、意谓的、主观的崇高内在性（书之尽）。书模仿世界，正如艺术模仿自然；运用其自身所特有的手段，产生出那些自然不能或不再能创造出来的东西。书的法则，则是反思／映像（réflexion）的法则，即“一”生成为“二”。那书的法则又怎能存在于自然之中，既然它统辖着世界与书、自然与艺术之间的区分？一生成二；每当我们遇到这个法则——即便被毛泽东策略性地提出，或被人们最为“辩证地”理解，我们所面临着的正是最为经典、最具反思性、最古老、最令人厌倦的思想。自然并不是这样运作的；在自然之中，根是直根，具有更为大量的侧面的或环状的（而非二元分化的）分支。精神落后于自然。甚至作为自然实在的书也是一种直根，有着垂直的轴和环绕的叶。然而书作为精神实在——树或根的形象——却不断展现着一生二、二生四的法则……二元逻辑是“树—根”的精神实在。即使像语言学这样“领先的”学科也仍然恪守着“树—根”这个基本形象，由此与反思的古典传统联结在一起（比如乔姆斯基和他的意群树：从一个点S开始，以二元分化的方式衍生）。可以说，此种思想从未理解多元性：它必需预设一种根本的、强有力的统一，以便遵循一种精神的方法来达到“二”。从客体的角度来说，遵循自然的方法，人们无疑可以直接从“一”达到三、四或五，但却始终必须预设一种根本的、强有力的统一性作为前提，即支撑着次级根的直根的统一性。这并不能使我们走得更远。取代二元分化之二元逻辑的，只不过是连续的循环之间的——对应的关系。无论是直根还是二歧根，都不能更好地解释多元性。前者运作于客体之中，后者运作于主体之中。二元逻辑和一一对应关系仍然统治着精神分析［弗洛伊德在解释施列伯（Schreber）病例时所运用的谵妄之树］、语言学、结构主义，乃至信息论。



胚根系统，或束根，是书的第二种形象，我们的现代性显然仰赖于它。这回，主根已然夭折，或者，它的末端已然被摧毁，一个任意的、直接的次级根的多元体被嫁接于它之上，并呈现出一种蓬勃的生长态势。这回，自然实在显现于主根夭折的过程中，但根的统一性仍然持有，作为过去或将来，作为可能性。我们须追问，是否反思性的精神实在没有以这样的方式对此种事态进行弥补：即转而诉求一种更为全面的隐秘统一性或更为广泛的总体性。不妨采用威廉·巴勒斯的剪裁法（cut-up）：将一篇文本叠合（pliage）进另一篇之中，由此构成了多元的，甚至是偶生的根（比如一根插穗），这对于所考察的文本来说就意味着一种替补的维度（dimension supplémentaire）。正是在这个叠合的替补维度之中，统一性继续着其精神的劳作。正是在这个意义上，最为碎片化的著作也同样可以被视为“全集”或“巨著”。绝大多数用来令系列得以增殖或令多元性得以拓展的现代方法在某个方向上（比如线性的方向）是极为有效的，而一种总体化的统一性却在另一个维度（循环或周期）之中获得了更为有力的肯定。每当一个多元体被掌控于一种结构之中时，其增长就被后者的组合法则所缩减甚至抵消。主张摧毁统一性的人也正是天使的缔造者，**天使博士（*doctores angelici*）\textbf{，因为他们肯定了一种真正的、至上的、天使般的统一性。乔伊斯的词语可谓名副其实地“具有多重根”，它们有效地粉碎了词语乃至语言的线性统一，但由此却得出了一种句子、文本或知识的循环统一。尼采的格言粉碎了知识的线性统一，但却由此回归于永恒轮回的循环统一［作为思想之中的非—知（non-su）］。如此说来，束根系统尚未真正摆脱二元论，它展现出了主客之间、自然实在和精神实在之间的互补性：在客体之中，统一性不断遭到阻碍甚至阻止，但在主体之中，一种新的统一性却获得胜利。世界失去了其主根，而主体也不再能够进行二元分化，但却在于其客体的维度始终构成替补的维度中，达到了一种更高的统一性，一种矛盾的或超定（surdétermination）的统一性。世界生成为混沌，但书仍然是世界的形象；胚根—混沌体（chaosmos），而不再是根—宇宙（cosmos）。说异的迷思：因其碎片化而更具有总体性的书。书作为世界的形象，无论如何都是乏味的观念。确实，仅仅说“多元性万岁”，这还是不够的，尽管要想发出这样的呼声也绝非易事。没有哪种排版、词语、甚或句法的技巧足以使其被听见。}必须形成“多”，但不是通过始终增加一个更高的维度，而是相反，以最为简单而节制的方式，始终对已有的维度进行 n-1 的操作（只有如此，一才成为多的构成部分，即始终是被减去）。从有待构成的多元体中减去独一无二者，在 n-1 的维度上写作。这样的体系可以被称为根茎（rhizome）。作为地下的茎，根茎截然有别于根和胚根。球茎和块茎都属于根茎。具有根或胚根的植物从所有其他的方面来看也可以是根茎式的；问题就在于，是否植物学从其特性上来说是完全是根茎式的。甚至某些动物也是根茎式的，在其集群（meute）的形态之中。鼠群就是根茎。兽穴也是根茎，它们有着各种栖居、储藏、移动、躲避、断裂的功能。根茎自身具有异常多样的形态，从表面向四面八方延展，直到凝聚成球茎和块茎的形态。当鼠群之间彼此窜动之时。存在着最好的和最糟的根茎：土豆，茅草和莠草。动物和植物，茅草就是螃蟹草（crab-grass）。我们觉得，如果不列举一些根茎的大致特征的话，是无法令人信服的。

\newpage
\vspace*{1em}
\begin{center}
{\Large\bfseries 课程综述}
\end{center}
\vspace{1em}
\addcontentsline{toc}{section}{课程综述}

从哪里开始，又走到了哪里。



\textbf{第一天。} 你们变成了一支田野调查队的成员。不是来评论一本小说，而是潜入陌生的权力场域。K 已经在村子里了，给你们写了第一封情报。你们接触了田野伦理，读了人类学调查方法的清单。你们意识到，田野调查是一种既疏离又热烈的思想姿态。观察与干预相伴而行，任何一次策略性抉择，都意味着真实之树的一次分叉。你们向城堡投去了最初的一瞥。



\textbf{第二天。} K 泊入田野的行动相当激进，旋即与田野中人建立亲密关系。这引发了你们的担忧和争论，是深入田野的重大进展，还是一步险棋。你们注意到 K 把同一个名字冠在两个老助手头上，一个叫阿图尔一个叫耶利米，但他们都叫“助手”。这是被动之举还是防范措施？你们感到不适，心怀求知的渴望却迷失于细节和疑团里，是煎熬的。但你们意识到，这正是调查的常态，关键在于自觉地意识到它。想要解释职能与名称的分离，就要理解政治分工有别于经济分工的奥秘。



\textbf{第三天。} K 完成了田野角色的确立。你们追问，远离权力中心，为什么反而是介入组织的更好途径。成为组织的一员，将对个体产生哪些心理后果。研究城堡、组织、权力，将如何有助于改善我们的社会生活。你们积累了很多问题，准备在后面的时间里去深化和解决它们。



\textbf{第四天。} 你们读了周雪光，了解了制度化了的政府间共谋现象。为什么组织权威默许共谋？讨论里出现了不同的声音，有人认为是有心无力鞭长莫及，有人认为相安无事就是更好的选择。你们理解了决策-执行的二元关系，执行对决策的扭曲不可避免，决策者必须在决策环节预判和吸纳这种扭曲。这才是非正式制度的真正来源。但你们没有停在这里，如果非正式的执行行为已经制度化，为什么不直接正式化？问题背后是体制之辩，是治理空间结构的潜在多元性。



\textbf{第五天。} 你们和纪君班合班上课，经历了一场从现象到理论的思想历程。你们分享了对教师角色的理解，意识到教师身处复杂的关系网络中，在制度之下努力维系着诸多关乎个人的良好意愿。但个体的良好意愿不易上升为组织行为。组织对确定性的偏爱促使教师倡导标准和可预测，哪怕他们把学生的差异和潜力看在心里。老师是在学生生活中栖居的街头官僚，正式职责之外，固有着一系列由栖居于他人生活而带来的非正式职责。两种职责并不平衡，这一定程度上源于家长特别的中介地位。你们也读了 1985 年的教育体制改革决定，教育被设定为服务于经济和行政建制的人才供应线。你们有理由相信，这样富有时代色彩的选择将随时间而改变。



\textbf{第六天。} 你们回到了城堡本身。不是在外部看，而是进入它的内部运作。电话是八音盒，打到城堡的电话让最下一级部门全部响起来，但接起来的人不一定是你找的人。公文在每个环节都有回应，但从未有裁断。下级负责遗忘，把决定权推出去，上级负责铭记，不断追问、催促。这组织的动态平衡。



\textbf{第七天。} 德勒兹与加塔利说，世界上有三种生命空间形态。树根，统一的生命体，从唯一起点出发，通过二元分化不断生长，对应的是韦伯式的正式组织。胚根，矛盾的生命体，多个起点相互嫁接，任何路径都可以被替代，整体在崩溃和重建之间轮回，对应的是列宁式的革命组织。根茎，多元的生命体，脱离起源，不断迁移、断裂，在任何一个节点上都蕴含着重新开始的潜力，对应的是卡夫卡式的非正式组织。



三种形态是理想型。真实世界中的组织从未只是其中一个。它们声称为韦伯式的正式组织，渴望成为列宁式的革命组织，而实际上是卡夫卡式的非正式组织。后者如幽灵般弥漫，难以捉摸，又如真相般难以抗拒。



这就是你们走过的七天。从田野调查者到组织观察者，从读一本书到理解一种存在方式。七天前你们打开《城堡》，可能觉得它是一个关于官僚制的寓言。现在你们知道，它不是寓言。它是田野笔记。



每个人看到的城堡都不一样。不是因为有人看对了有人看错了，而是因为理解组织就是理解自己理解它的方式。你们意识到这一点，这七天就值得了。



谢谢你们的思想劳作，感谢这一程的陪伴！

\end{document}